% ID: SEL-D0058-OA-E01
% ESTADO: cotejada. Segunda pasada visual y borrador compilado revisados.
% Fuente: Selectividad/sel_2011_matematicas_aplicadas/221A-EXAMEN 5.pdf
% SHA-256: 1780d560801fb3001e146f12b3a37b29b566ca6576dc9cdafda07902d7dee982
% Página: 1; opción A; ejercicio original 1.

Se considera el recinto $R$ del plano, determinado por las siguientes inecuaciones:
\[
x+y\geq2,\quad x+3y\leq15,\quad3x-y\leq15,\quad x\geq0,\quad y\geq0.
\]
\begin{enumerate}[label=\alph*)]
\item \textbf{(1.5 puntos)} Represente gráficamente el recinto $R$ y calcule sus vértices.
\item \textbf{(0.5 puntos)} Halle los valores máximo y mínimo que alcanza la función $F(x,y)=3x+y$ en dicho recinto.
\item \textbf{(0.5 puntos)} Razone si existen puntos $(x,y)$ del recinto, para los que $F(x,y)=30$.
\end{enumerate}
